\section{Tesis 4}

\subsection{Título}
\textit{Diseño e implementación de un método para la ejecución de aplicaciones de computación paralela basadas en MPI en un clúster de Kubernetes} \cite{arias2024}.

\subsection{Autor}
Cristian Arias Chaves, Universidad de Costa Rica.

\subsection{Resumen}
La tesis diseña, implementa y valida un método para ejecutar aplicaciones de computación paralela basadas en MPI (\textit{Message Passing Interface}) dentro de un clúster de Kubernetes. El método resuelve los requisitos técnicos de MPI en un entorno contenedorizado (IPs y DNS estables, almacenamiento compartido, ambiente de ejecución uniforme) mediante una combinación de recursos nativos de Kubernetes: \textit{StatefulSets}, servicios \textit{headless}, volúmenes persistentes con BeeGFS e imágenes de contenedor auto-contenidas. El método es automatizado mediante un Operador de Kubernetes desarrollado en Go usando el \textit{Operator SDK}, con dos Definiciones de Recursos Personalizados (CRDs): \texttt{MpiCluster} y \texttt{MpiJob}. La validación se realiza mediante cuatro experimentos: prueba funcional básica (``Hola Mundo'' de MPI), medición del tiempo de creación del clúster, medición del rendimiento de entrada/salida de red y disco, y evaluación del rendimiento de ejecución de aplicaciones con los \textit{benchmarks} MiniFE, MiniMD y MiniXyce de la suite Mantevo. Los tres tipos de imágenes de contenedor probadas (MPICH, OpenMPI y MPICH+SLURM) demuestran que el clúster de trabajo en Kubernetes puede crearse en aproximadamente 21 segundos por nodo, con rendimientos de E/S y de ejecución de aplicaciones superiores al 70\% del rendimiento de un clúster físico equivalente en todos los casos \cite{arias2024}.

\subsection{Problema que busca resolver}
Los clústeres HPC (\textit{High Performance Computing}) tradicionales son infraestructura cara, rígida y especializada que requiere inversión de capital significativa y personal técnico dedicado. Kubernetes se ha convertido en el estándar de orquestación de contenedores en la nube, pero no fue diseñado para los requisitos específicos de la computación paralela con MPI: los \textit{deployments} de Kubernetes no garantizan IPs ni DNS estables entre nodos (requisito de MPI para el intercambio de mensajes), no proveen almacenamiento compartido de forma nativa, y la configuración del ambiente de ejecución para cada aplicación paralela debe realizarse manualmente. No existe un método estandarizado, reproducible y productivo para desplegar y gestionar trabajos MPI en Kubernetes, lo que crea una barrera entre el ecosistema moderno de la nube y la comunidad de computación de alto rendimiento.

\subsection{Objetivos}

Objetivo General: Diseñar e implementar un método que permita la ejecución de aplicaciones de computación paralela basadas en MPI en un clúster de Kubernetes, con un rendimiento de E/S y de ejecución de aplicaciones no inferior al 70\% respecto a un clúster físico equivalente, y con un tiempo de creación e inicialización de cada nodo del clúster de trabajo no superior a 1 minuto \cite{arias2024}.

Objetivos Específicos:
\begin{enumerate}
    \item Identificar los requisitos técnicos para la correcta ejecución de aplicaciones MPI en un sistema distribuido y determinar los recursos de Kubernetes más adecuados para satisfacerlos.
    \item Diseñar e implementar el método propuesto usando recursos nativos de Kubernetes (\textit{StatefulSet}, servicio \textit{headless}, volumen persistente, \textit{namespace}, cuotas de recursos) e imágenes de contenedor auto-contenidas.
    \item Validar el método en un ambiente de desarrollo demostrando la ejecución exitosa de la implementación de Hola Mundo de MPI en todos los nodos del clúster de trabajo.
    \item Medir el tiempo de creación del clúster de trabajo y verificar que no supere 1 minuto por nodo ni $1 \times N$ minutos para $N$ nodos.
    \item Evaluar el rendimiento de E/S de disco y red del clúster de trabajo en comparación con un clúster físico de Vultr y verificar que supera el 70\%.
    \item Evaluar el rendimiento de la ejecución de aplicaciones de computación paralela mediante \textit{benchmarks} de la suite Mantevo y verificar que supera el 70\% del clúster físico.
\end{enumerate}

\subsection{Hipótesis}

La tesis formula explícitamente tres hipótesis cuantitativas, lo que permite una evaluación objetiva de los resultados:

\begin{enumerate}
    \item \textit{El tiempo de creación e inicialización de cada nodo del clúster de trabajo no superará 1 minuto (60 segundos), y el tiempo total para iniciar la ejecución de una aplicación de computación paralela será igual o menor a $1 \times N$ minutos, donde $N$ es el número de nodos del clúster de trabajo.} Confirmada: el tiempo promedio fue de aproximadamente 21 segundos por nodo para las tres imágenes evaluadas.

    \item \textit{El rendimiento de entrada/salida de disco y red en los nodos del clúster de trabajo de Kubernetes superará el 70\% del rendimiento alcanzado en los nodos del clúster físico de Vultr.} Confirmada: el rendimiento de disco fue siempre superior al 70\% (la mayoría de nodos superó el 90\%) y el rendimiento de red superó incluso el 100\% del clúster físico en todos los nodos.

    \item \textit{El rendimiento de ejecución de aplicaciones de computación paralela basadas en MPI en el clúster de trabajo de Kubernetes superará el 70\% del rendimiento alcanzado en el clúster físico de Vultr.} Confirmada: todos los benchmarks alcanzaron rendimientos promedios superiores al 84\% en escalamiento fuerte y al 91\% en escalamiento débil.
\end{enumerate}

Las tres hipótesis habrían sido refutadas si los tiempos de creación hubieran superado 60 segundos, si el rendimiento de E/S o de ejecución hubiera caído por debajo del 70\%, o si el sistema hubiera fallado en ejecutar correctamente las aplicaciones MPI \cite{arias2024}.

\subsection{Resumen de la metodología empleada}

La metodología empleada por \cite{arias2024} se describe a continuación:

\begin{itemize}
    \item \textbf{Diseño del método:} Para resolver los requisitos técnicos de la computación paralela dentro de Kubernetes, se usaron recursos especializados de la plataforma que garantizan que cada servidor tenga una dirección de red estable y un nombre fijo, necesario para que los procesos paralelos puedan comunicarse entre sí. El almacenamiento compartido entre servidores se implementó con un sistema de archivos de alto rendimiento diseñado para computación paralela (BeeGFS). Cada servidor ejecuta un contenedor auto-contenido que incluye el sistema operativo y las librerías necesarias para correr los programas paralelos.

    \item \textbf{Automatización mediante un programa gestor:} Se desarrolló un programa de automatización que gestiona todo el ciclo de vida del clúster: crearlo, configurarlo y eliminarlo. Este programa permite al usuario definir la configuración del clúster (cantidad de servidores, tipo de software y almacenamiento) y los trabajos a ejecutar, sin necesidad de intervención manual en cada paso.

    \item \textbf{Configuraciones de software evaluadas:} Se probaron tres variantes de software para ejecutar programas paralelos: dos implementaciones estándar del protocolo de comunicación entre procesos (MPICH y OpenMPI), y una tercera que además incluye un gestor de colas de trabajos para entornos más complejos.

    \item \textbf{Infraestructura:} El clúster en Kubernetes se montó en la nube sobre ocho servidores de 15 núcleos y 30 GB de memoria RAM cada uno. Como referencia de comparación se usó el mismo proveedor de nube pero con los mismos servidores configurados sin Kubernetes, para simular un clúster físico equivalente.

    \item \textbf{Experimentos de validación:}
    \begin{enumerate}
        \item Prueba funcional básica: se verificó que el clúster ejecutara correctamente un programa paralelo sencillo en todos los servidores.
        \item Tiempo de arranque: se midió cuántos segundos tarda en quedar listo cada servidor del clúster.
        \item Velocidad de disco: se midió la velocidad de lectura y escritura de archivos en disco en cada servidor.
        \item Velocidad de red: se midió la velocidad de transferencia de datos entre servidores.
        \item Rendimiento de aplicaciones: se ejecutaron tres programas de simulación científica variando la cantidad de procesos paralelos activos (de 1 a 120), para comparar el rendimiento en Kubernetes contra el del servidor físico de referencia.
    \end{enumerate}
\end{itemize}

\subsection{Principales resultados}

Los principales hallazgos reportados por \cite{arias2024} son:

\begin{itemize}
    \item \textbf{Tiempo de arranque del clúster:} Las tres configuraciones de software probadas lograron tiempos promedio de aproximadamente 21 segundos por servidor, muy por debajo del límite de 60 segundos establecido como criterio de éxito. El primer servidor tardó un poco más (entre 26 y 29 segundos) porque es el encargado de generar las credenciales de seguridad para todos los demás nodos del clúster.

    \item \textbf{Velocidad de lectura y escritura en disco:} El rendimiento del clúster en Kubernetes superó el 70\% del rendimiento del servidor físico de referencia en todos los nodos; la mayoría alcanzó más del 90\%. Los resultados fueron consistentes entre repeticiones, lo que indica estabilidad del sistema.

    \item \textbf{Velocidad de transferencia de datos por red:} El clúster en Kubernetes superó al servidor físico de referencia en todos los nodos, con rendimientos de entre el 102\% y el 135\%. Esto se atribuye a que las pruebas del servidor físico posiblemente se realizaron en momentos de mayor congestión de red, o a que el componente de red de Kubernetes transmite datos de forma más eficiente en esa configuración.

    \item \textbf{Rendimiento ejecutando simulaciones --- MiniFE:} Con un programa de simulación de elementos finitos, el clúster en Kubernetes alcanzó un rendimiento promedio del 97.48\% respecto al servidor físico cuando el tamaño del problema era fijo y se aumentaban los procesos paralelos, y del 91.01\% cuando el tamaño del problema crecía junto con los procesos. En ambos casos el comportamiento fue casi idéntico al del servidor físico.

    \item \textbf{Rendimiento ejecutando simulaciones --- MiniMD:} Con un simulador de dinámica molecular, el promedio fue del 84.76\% con problema de tamaño fijo y del 97.58\% con tamaño proporcional. El único resultado por debajo del 70\% en todo el estudio fue un caso aislado con 45 procesos paralelos en esta prueba (65.3\%), que la tesis reconoce como atípico sin ofrecer una explicación técnica definitiva.
\end{itemize}

\subsection{Crítica de la tesis}
\begin{itemize}
    \item \textbf{Aspectos positivos:} Las secciones siguen una secuencia lógica clara de diseño, implementación y validación, coherente con el tipo de tesis de ingeniería aplicada. Las tres hipótesis son cuantitativas y verificables, se corresponden directamente con los objetivos específicos y las conclusiones las confirman de forma explícita, cerrando el ciclo del trabajo con rigor. El título describe con precisión el alcance del método propuesto y el alcance es consistente con el trabajo relacionado: la tesis posiciona correctamente la brecha entre el HPC tradicional y el ecosistema Kubernetes antes de proponer su solución, lo que justifica la pertinencia del trabajo dentro de la literatura existente.
    \item \textbf{Aspectos de posible mejora:} En primer lugar, la comparación se realiza únicamente contra el mismo proveedor de nube sin Kubernetes, sin contrastar el método frente a soluciones alternativas existentes como el \textit{Kubeflow MPI Operator} o Slurm; esto limita el posicionamiento del trabajo en el estado del arte y debilita las conclusiones respecto a su ventaja competitiva. En segundo lugar, el criterio de éxito del 70\% de rendimiento no se fundamenta con referencias a estándares del área de HPC en la nube, por lo que no queda claro si el umbral es exigente o conservador dentro de la literatura. Por último, el resultado atípico de MiniMD con 45 \textit{ranks} (65.3\%, único valor inferior al 70\%) no recibe una explicación técnica satisfactoria, lo que genera una duda sobre la robustez del método ante ciertas configuraciones de nodos.
\end{itemize}
